\chapter{2008年四川延考卷（文）}
\section{单选题}

\begin{problem}
设集合 \( A = \{-1, 0, 1\} \)， \(A\) 的子集中，含有元素 \(0\) 的子集共有（\quad）

\choices
    {\(2\) 个}
    {\(4\) 个}
    {\(6\) 个}
    {\(9\) 个}
\end{problem}

\begin{answer}
B
\end{answer}

\begin{solution}
要使子集含有元素 \(0\)，元素 \(0\) 必须选入；另外两个元素 \(-1\)，\(1\) 可以分别选入或不选入，因此共有 \(2^2=4\) 个这样的子集，分别为 \(\{0\}\)，\(\{-1,0\}\)，\(\{0,1\}\)，\(\{-1,0,1\}\). 故选 B.
\end{solution}

\begin{problem}
函数 \( y=\sqrt{1-x}+\lg x \) 的定义域为（\quad）

\choices
    {\( (0,+\infty) \)}
    {\( (-\infty,1] \)}
    {\( (-\infty,0)\cup[1,+\infty) \)}
    {\( (0,1] \)}
\end{problem}

\begin{answer}
D
\end{answer}

\begin{solution}
平方根有意义要求 \(1-x\geq 0\)，即 \(x\leq 1\)；对数 \(\lg x\) 有意义要求 \(x>0\). 两个条件取交集，得定义域为 \((0,1]\). 故选 D.
\end{solution}

\begin{problem}
\( \left(1+\frac{1}{x}\right)\left(1+x\right)^{4} \) 的展开式中含 \( x^{2} \) 项的系数（\quad）

\choices
    {\(4\)}
    {\(6\)}
    {\(10\)}
    {\(12\)}
\end{problem}

\begin{answer}
C
\end{answer}

\begin{solution}
展开式中 \(x^2\) 项有两种来源：从 \((1+x)^4\) 中取 \(x^2\) 项，系数为 \(\mathrm C_4^2=6\)；从 \(\dfrac{1}{x}(1+x)^4\) 中取 \(x^2\) 项，须从 \((1+x)^4\) 中取 \(x^3\) 项，系数为 \(\mathrm C_4^3=4\). 所以 \(x^2\) 项的系数为 \(6+4=10\). 故选 C.
\end{solution}

\begin{problem}
不等式 \( |x-2|<1 \) 的解集为（\quad）

\choices
    {\( \left\{x \mid 1<x<3\right\} \)}
    {\( \left\{x \mid 0<x<2\right\} \)}
    {\( \left\{x \mid 1<x<2\right\} \)}
    {\( \left\{x \mid 2<x<3\right\} \)}
\end{problem}

\begin{answer}
A
\end{answer}

\begin{solution}
由绝对值不等式的定义，\(|x-2|<1\) 等价于 \(-1<x-2<1\). 不等式两边同加 \(2\)，得 \(1<x<3\)，所以解集为 \(\{x\mid 1<x<3\}\). 故选 A.
\end{solution}

\begin{problem}
已知 \( \tan\alpha=\frac{1}{2} \)，则 \( \frac{(\sin\alpha+\cos\alpha)^{2}}{\cos2\alpha}= \)（\quad）

\choices
    {\(2\)}
    {\(-2\)}
    {\(3\)}
    {\(-3\)}
\end{problem}

\begin{answer}
C
\end{answer}

\begin{solution}
由 \(\tan\alpha=\dfrac{1}{2}\)，且正切有意义，知 \(\cos\alpha\ne0\). 因而
\[
\frac{(\sin\alpha+\cos\alpha)^2}{\cos2\alpha}=\frac{\cos^2\alpha(\tan\alpha+1)^2}{\cos^2\alpha-\sin^2\alpha}=\frac{\left(\frac12+1\right)^2}{1-\left(\frac12\right)^2}=\frac{9}{4}\mathbin{/}\frac{3}{4}=3
\]
故选 C.
\end{solution}

\begin{problem}
一个三棱锥的底面边长等于一个球的半径，该正三棱锥的高等于是这个球的直径，则球的体积与正三棱锥的体积的比值为（\quad）

\choices
    {\( \frac{8\sqrt{3}\pi}{3} \)}
    {\( \frac{\sqrt{3}\pi}{6} \)}
    {\( \frac{\sqrt{3}\pi}{2} \)}
    {\( 8\sqrt{3}\pi \)}
\end{problem}

\begin{answer}
A
\end{answer}

\begin{solution}
设球的半径为 \(r\)，球的体积记为 \(V_1\)，正三棱锥的体积记为 \(V_2\). 球的体积为 \(V_1=\dfrac{4}{3}\pi r^3\). 正三棱锥的底面是边长为 \(r\) 的正三角形，底面积为 \(\dfrac{\sqrt3}{4}r^2\)，高为球的直径 \(2r\)，所以其体积为
\[
V_2=\frac13\cdot\frac{\sqrt3}{4}r^2\cdot2r=\frac{\sqrt3}{6}r^3
\]
于是
\[
\frac{V_1}{V_2}=\frac{\frac43\pi r^3}{\frac{\sqrt3}{6}r^3}=\frac{8\sqrt3\pi}{3}
\]
故选 A.
\end{solution}

\begin{problem}
若点 \(P(2,0)\) 到双曲线 \( \frac{x^{2}}{a^{2}} - \frac{y^{2}}{b^{2}} = 1 \) 的一条渐近线的距离为 \( \sqrt{2} \)，则双曲线的离心率为（\quad）

\choices
    {\( \sqrt{2} \)}
    {\( \sqrt{3} \)}
    {\( 2\sqrt{2} \)}
    {\( 2\sqrt{3} \)}
\end{problem}

\begin{answer}
A
\end{answer}

\begin{solution}
双曲线的一条渐近线可写为 \(y=\dfrac{b}{a}x\)，即 \(bx-ay=0\)，其中 \(a>0\)，\(b>0\). 点 \(P(2,0)\) 到该直线的距离为
\[
\frac{|2b|}{\sqrt{a^2+b^2}}=\sqrt2
\]
两边平方得 \(\dfrac{4b^2}{a^2+b^2}=2\)，即 \(b^2=a^2\). 双曲线的离心率满足
\[
e=\frac{c}{a}=\sqrt{1+\frac{b^2}{a^2}}=\sqrt2
\]
故选 A.
\end{solution}

\begin{problem}
在一次读书活动中，一个同学从 \(4\) 本不同的科技书和 \(2\) 本不同的文艺书中任选 \(3\) 本，则所选的书中既有科技书又有文艺书的概率为（\quad）

\choices
    {\( \frac{1}{5} \)}
    {\( \frac{1}{2} \)}
    {\( \frac{2}{3} \)}
    {\( \frac{4}{5} \)}
\end{problem}

\begin{answer}
D
\end{answer}

\begin{solution}
从 \(6\) 本不同的书中任选 \(3\) 本，共有 \(\mathrm C_6^3=20\) 种选法. 既有科技书又有文艺书时，分为选 \(2\) 本科技书和 \(1\) 本文艺书，或选 \(1\) 本科技书和 \(2\) 本文艺书两类，故有
\[
\mathrm C_4^2\mathrm C_2^1+\mathrm C_4^1\mathrm C_2^2=12+4=16
\]
种. 所求概率为 \(\dfrac{16}{20}=\dfrac45\). 故选 D.
\end{solution}

\begin{problem}
过点 \((0,1)\) 的直线与圆 \(x^{2} + y^{2} = 4\) 相交于 \(A\)，\(B\) 两点，则 \(|AB|\) 的最小值为（\quad）

\choices
    {\(2\)}
    {\(2\sqrt{3}\)}
    {\(3\)}
    {\(2\sqrt{5}\)}
\end{problem}

\begin{answer}
B
\end{answer}

\begin{solution}
圆心为原点，半径为 \(2\). 过点 \(P(0,1)\) 的任意直线到圆心的距离 \(d\) 满足 \(0\leq d\leq OP=1\)，且当直线垂直于 \(OP\) 时取最大值 \(d=1\). 该直线截圆所得弦长为
\[
|AB|=2\sqrt{2^2-d^2}
\]
要使弦长最小，应使 \(d\) 最大，因此
\[
|AB|_{\min}=2\sqrt{4-1}=2\sqrt3
\]
故选 B.
\end{solution}

\begin{problem}
已知两个单位向量 \(\bs{a}\) 与 \(\bs{b}\) 的夹角为 \( \frac{\pi}{3} \)， \(\bs{a} + \lambda \bs{b} \) 与 \(\lambda \bs{a} - \bs{b} \) 互相垂直的充要条件是（\quad）

\choices
    {\( \lambda = -\frac{\sqrt{3}}{2} \) 或 \( \lambda = \frac{\sqrt{3}}{2} \)}
    {\( \lambda = -\frac{1}{2} \) 或 \( \lambda = \frac{1}{2} \)}
    {\( \lambda = -1 \) 或 \( \lambda = 1 \)}
    {\( \lambda \) 为任意实数}
\end{problem}

\begin{answer}
C
\end{answer}

\begin{solution}
由 \(|\bs a|=|\bs b|=1\) 及夹角为 \(\dfrac\pi3\)，得 \(\bs a\cdot\bs b=\cos\dfrac\pi3=\dfrac12\). 两向量互相垂直的充要条件是数量积为零，所以
\[
(\bs a+\lambda\bs b)\cdot(\lambda\bs a-\bs b)=\lambda-\frac12+\frac{\lambda^2}{2}-\lambda=\frac{\lambda^2-1}{2}=0
\]
因此 \(\lambda=\pm1\)，即 \(\lambda=-1\) 或 \(\lambda=1\). 故选 C.
\end{solution}

\begin{problem}
设函数 \( y = f(x) \)（\( x \in \R \)）的图像关于直线 \(x = 0\) 及直线 \(x = 1\) 对称，且 \( x \in [0,1] \) 时， \( f(x) = x^2 \)，则 \( f\left(-\frac{3}{2}\right) = \)（\quad）

\choices
    {\( \frac{1}{2} \)}
    {\( \frac{1}{4} \)}
    {\( \frac{3}{4} \)}
    {\( \frac{9}{4} \)}
\end{problem}

\begin{answer}
B
\end{answer}

\begin{solution}
关于直线 \(x=0\) 对称，故 \(f\left(-\dfrac32\right)=f\left(\dfrac32\right)\). 再关于直线 \(x=1\) 对称，点 \(\dfrac32\) 关于该直线的对称点为 \(\dfrac12\)，所以 \(f\left(\dfrac32\right)=f\left(\dfrac12\right)\). 由于 \(\dfrac12\in[0,1]\)，有
\[
f\left(-\frac32\right)=f\left(\frac12\right)=\left(\frac12\right)^2=\frac14
\]
故选 B.
\end{solution}

\begin{problem}
在正方体 \( ABCD-A_{1}B_{1}C_{1}D_{1} \) 中，\(E\) 是棱 \(A_{1}B_{1}\) 的中点，则 \(A_{1}B\) 与 \(D_{1}E\) 所成角的余弦值为（\quad）

\choices
    {\( \frac{\sqrt{5}}{10} \)}
    {\( \frac{\sqrt{10}}{10} \)}
    {\( \frac{\sqrt{5}}{5} \)}
    {\( \frac{\sqrt{10}}{5} \)}
\end{problem}

\begin{answer}
B
\end{answer}

\begin{solution}
设正方体棱长为 \(l\)，建立坐标系
取 \(A(0,0,0)\)，\(B(l,0,0)\)，\(A_1(0,0,l)\)，\(D_1(0,l,l)\)，\(E\left(\frac l2,0,l\right)\).
于是 \(\overrightarrow{A_1B}=(l,0,-l)\)，\(\overrightarrow{D_1E}=\left(\frac l2,-l,0\right)\).
它们的数量积为 \(\dfrac{l^2}{2}\)，长度分别为 \(l\sqrt2\) 和 \(\dfrac{l\sqrt5}{2}\). 因而两直线所成角的余弦值为
\[
\frac{\left|\overrightarrow{A_1B}\cdot\overrightarrow{D_1E}\right|}{|\overrightarrow{A_1B}|\,|\overrightarrow{D_1E}|}=\frac{\frac{l^2}{2}}{l\sqrt2\cdot\frac{l\sqrt5}{2}}=\frac{\sqrt{10}}{10}
\]
故选 B.
\end{solution}

\section{填空题}

\begin{problem}
函数 \( y = \e^{x+1} - 1 \)（\( x \in \R \)）的反函数为\fillinblank{}.
\end{problem}

\begin{answer}
\( y=\ln(x+1)-1 \)（\(x>-1\)）
\end{answer}

\begin{solution}
由 \(y=\e^{x+1}-1\)，得 \(y+1=\e^{x+1}\)，所以 \(x=\ln(y+1)-1\). 交换 \(x\)，\(y\) 得反函数
\[
y=\ln(x+1)-1
\]
原函数的值域由 \(\e^{x+1}>0\) 得为 \((-1,+\infty)\)，故反函数的定义域为 \(x>-1\).
\end{solution}

\begin{problem}
函数 \( f(x)=\sqrt{3}\sin x-\cos^{2}x \) 的最大值是\fillinblank{}.
\end{problem}

\begin{answer}
\( \sqrt{3} \)
\end{answer}

\begin{solution}
令 \(t=\sin x\)，则 \(t\in[-1,1]\)，且 \(\cos^2x=1-t^2\). 因而
\[
f(x)=\sqrt3t-(1-t^2)=t^2+\sqrt3t-1
\]
这是关于 \(t\) 的开口向上的二次函数，其对称轴为 \(t=-\dfrac{\sqrt3}{2}\)，所以在区间 \([-1,1]\) 上的最大值出现在端点. 计算得
\(\left.t^2+\sqrt3t-1\right|_{t=1}=\sqrt3\)，\(\left.t^2+\sqrt3t-1\right|_{t=-1}=-\sqrt3\).
故函数的最大值为 \(\sqrt3\).
\end{solution}

\begin{problem}
设等差数列 \( \left\{a_{n}\right\} \) 的前 \(n\) 项和为 \( S_{n} \)，且 \( S_{5}=a_{5} \). 若 \( a_{4}\neq0 \)，则 \( \frac{a_{7}}{a_{4}}= \)\fillinblank{}.
\end{problem}

\begin{answer}
\(3\)
\end{answer}

\begin{solution}
设等差数列的首项为 \(a_1\)，公差为 \(d\). 则 \(S_5=\frac52(2a_1+4d)=5a_1+10d\)，\(a_5=a_1+4d\).
由 \(S_5=a_5\) 得 \(4a_1+6d=0\)，即 \(a_1=-\dfrac32d\). 因而
\(a_4=a_1+3d=\frac32d\)，\(a_7=a_1+6d=\frac92d\).
题设 \(a_4\ne0\) 保证 \(d\ne0\)，所以
\[
\frac{a_7}{a_4}=\frac{\frac92d}{\frac32d}=3
\]
\end{solution}

\begin{problem}
已知 \( \angle AOB = 90^\circ \)， \(C\) 为空间中一点，且 \( \angle AOC = \angle BOC = 60^\circ \). 则直线 \(OC\) 与 \(AOB\) 平面所成角的正弦值为\fillinblank{}.
\end{problem}

\begin{answer}
\( \frac{\sqrt{2}}{2} \)
\end{answer}

\begin{solution}
取单位向量 \(\bs u\)，\(\bs v\) 分别沿 \(OA\)，\(OB\) 方向. 因为 \(\angle AOB=90^\circ\)，所以 \(\bs u\)，\(\bs v\) 是平面 \(AOB\) 内的互相垂直单位向量. 设 \(\bs w\) 为沿 \(OC\) 方向的单位向量，有 \(\bs w\cdot\bs u=\cos60^\circ=\frac12\)，\(\bs w\cdot\bs v=\cos60^\circ=\frac12\).
因此 \(\bs w\) 在平面 \(AOB\) 上的正投影长度的平方为
\[
\left(\frac12\right)^2+\left(\frac12\right)^2=\frac12
\]
由于 \(|\bs w|=1\)，其垂直于平面 \(AOB\) 的分量长度为 \(\sqrt{1-\frac12}=\dfrac{\sqrt2}{2}\). 直线与平面所成角的正弦等于单位方向向量的法向分量长度，故所求值为 \(\dfrac{\sqrt2}{2}\).
\end{solution}

\section{解答题}

\begin{problem}
在 \(\triangle ABC\) 中，内角 \(A\)，\(B\)，\(C\) 对边的边长分别是 \(a\)，\(b\)，\(c\)，已知 \(a^2 + b^2 = 2b^2\).

\begin{enumerate}
\item 若 \(B=\frac{\pi}{4}\)，且 \(A\) 为钝角，求内角 \(A\) 与 \(C\) 的大小；

\item 求 \(\sin B\) 的最大值.
\end{enumerate}
\end{problem}

\begin{solution}
\begin{enumerate}
\item 由题设 \(a^2+b^2=2b^2\)，得 \(a^2=b^2\). 因为三角形边长为正数，所以 \(a=b\). 由正弦定理，\(a=b\) 等价于 \(\sin A=\sin B\). 又 \(A+B<\pi\)，故不可能有 \(A=\pi-B\)，只能有 \(A=B\). 当 \(B=\dfrac\pi4\) 时得到 \(A=\dfrac\pi4\)，这与 \(A\) 为钝角矛盾. 所以按原 PDF 题面，本小问不存在满足条件的三角形.

\item 仍由 \(a=b\) 得 \(A=B\). 由三角形内角和，\(C=\pi-2B>0\)，所以 \(0<B<\dfrac\pi2\). 从而 \(\sin B<1\). 另一方面，对任意 \(B\in\left(0,\dfrac\pi2\right)\)，取 \(A=B\)，\(C=\pi-2B\)，并按正弦定理取边长比 \(a:b:c=\sin B:\sin B:\sin2B\)，即可得到满足 \(a=b\) 的非退化三角形. 令 \(B\to\dfrac\pi2^{-}\)，则 \(\sin B\to1\). 因此 \(\sin B\) 没有最大值，上确界为 \(1\).
\end{enumerate}
\end{solution}

\begin{problem}
一条生产线上生产的产品按质量情况分为三类：A 类，B 类，C 类. 检验员定时从该生产线上任取 \(2\) 件产品进行一次抽检，若发现其中含有 C 类产品或 \(2\) 件都是 B 类产品，就需要调整设备，否则不需要调整. 已知该生产线上生产的每件产品为 A 类品，B 类品和 C 类品的概率分别为 \(0.9\)，\(0.05\) 和 \(0.05\)，各件产品的质量情况互不影响.

\begin{enumerate}
\item 求在一次抽检后，设备不需要调整的概率；

\item 若检验员一天抽检 \(3\) 次，求一天中至少有一次需要调整设备的概率.
\end{enumerate}
\end{problem}

\begin{solution}
\begin{enumerate}
\item 设 \(A_i\) 表示第 \(i\) 件产品为 \(A\) 类品，\(B_i\) 表示第 \(i\) 件产品为 \(B\) 类品，其中 \(i=1\)，\(2\). 设事件 \(E\) 表示一次抽检后设备不需要调整. 由题意，只有抽到 \(AA\)，\(AB\) 或 \(BA\) 时不需要调整，因此
\[
P(E)=P(A_1\cap A_2)+P(A_1\cap B_2)+P(B_1\cap A_2)
=0.9^2+0.9\times0.05+0.05\times0.9=0.9
\]

\item 由（1）知，每次抽检不需要调整的概率为 \(0.9\). 各件产品质量情况相互独立，故三次抽检的结果也相互独立. 一天中至少有一次需要调整设备是“一天三次都不属于事件 \(E\)”的对立事件，所以所求概率为
\[
1-P(E)^3=1-0.9^3=1-0.729=0.271
\]
\end{enumerate}
\end{solution}

\begin{problem}
如图，一张平行四边形的硬纸片 \(ABC_0D\) 中，\(AD=BD=1\)，\(AB=\sqrt{2}\). 沿它的对角线 \(BD\) 把 \(\triangle BDC_0\) 折起，使点 \(C_0\) 到达平面 \(ABC_0D\) 外点 \(C\) 的位置.

\begin{enumerate}
\item 证明：平面 \(ABC_0D \perp\) 平面 \(CBC_0\).

\item 当二面角 \(A-BD-C\) 为 \(120^\circ\) 时，求 \(AC\) 的长.
\end{enumerate}

\bitmapfigure[width=0.58\linewidth]{img/2008/sichuan_postponed_liberal/q19_fig1.png}
\end{problem}

\begin{solution}
\begin{enumerate}
\item 证明：因为 \(AD=BC_0=BD=1\)，\(AB=C_0D=\sqrt{2}\)，所以 \(\angle DBC_0=90^\circ\)，\(\angle ADB=90^\circ\). 因为折叠过程中，\(\angle DBC=\angle DBC_0=90^\circ\)，所以 \(DB \perp BC\)，又 \(DB \perp BC_0\)，故 \(DB \perp\) 平面 \(CBC_0\). 又 \(DB \subset\) 平面 \(ABC_0D\)，所以平面 \(ABC_0D \perp\) 平面 \(CBC_0\).

\item 解法一：如图，由（1）知 \(BC \perp DB\)，\(BC_0 \perp DB\)，所以 \(\angle CBC_0\) 是二面角 \(C-BD-C_0\) 的平面角. 因为 \(A\) 与 \(C_0\) 在直线 \(BD\) 两侧，二面角 \(A-BD-C\) 与 \(C_0-BD-C\) 互为补角，故由已知 \(\angle A-BD-C=120^\circ\) 得 \(\angle CBC_0=60^\circ\). 作 \(CF \perp C_0B\)，垂足为 \(F\). 由 \(BC=BC_0=1\)，可得 \(CF=\frac{\sqrt{3}}{2}\)，\(BF=\frac{1}{2}\). 又由 \(\triangle ADB\) 为等腰直角三角形得 \(\angle ABD=45^\circ\)，结合 \(BD \perp BC_0\) 得 \(\angle ABF=\angle ABC_0=135^\circ\). 连结 \(AF\)，在 \(\triangle ABF\) 中，
\[
AF^2=\left(\sqrt{2}\right)^2+\left(\frac{1}{2}\right)^2-2\times \sqrt{2}\times \frac{1}{2}\times \cos 135^\circ=\frac{13}{4}
\]
因为平面 \(ABC_0D \perp\) 平面 \(CBC_0\)，所以 \(CF \perp\) 平面 \(ABC_0D\)，可知 \(CF \perp AF\). 在 \(\triangle AFC\) 中，
\[
AC=\sqrt{AF^2+CF^2}=\sqrt{\frac{13}{4}+\frac{3}{4}}=2
\]

解法二：由已知得 \(\angle ADB=\angle DBC_0=90^\circ\). 以 \(D\) 为坐标原点，射线 \(DA\)，\(DB\) 分别为 \(x\) 轴，\(y\) 轴的正半轴，建立如图的空间直角坐标系 \(D-xyz\). 则 \(A(1,0,0)\)，\(B(0,1,0)\)，\(C_0(-1,1,0)\)，\(D(0,0,0)\). 由（1）知 \(BC \perp DB\)，\(BC_0 \perp DB\)，所以 \(\angle CBC_0\) 是二面角 \(C-BD-C_0\) 的平面角；由上面的补角关系，\(\angle CBC_0=60^\circ\)，故 \(C\left(-\frac{1}{2},1,\frac{\sqrt{3}}{2}\right)\). 所以

\(\left| \overrightarrow{AC} \right|=\sqrt{\left(-\frac{1}{2}-1\right)^2+1^2+\left(\frac{\sqrt{3}}{2}\right)^2}=2\).
即 \(AC\) 的长为 \(2\).
\end{enumerate}
\FigureLayoutDeclare{img/2008/sichuan_postponed_liberal/q19_solution_fig1.png}{7.0cm}{26bp 19bp 69bp 19bp}
\solutionfigure{\bitmapfigure[width=0.50\linewidth]{img/2008/sichuan_postponed_liberal/q19_solution_fig1.png}}
\solutionfigure{\bitmapfigure[width=0.34\linewidth]{img/2008/sichuan_postponed_liberal/q19_solution_fig2.png}}
\end{solution}

\begin{problem}
在数列 \(\{a_n\}\) 中，\(a_1=1\)，\(2a_{n+1}=\left(1+\frac{1}{n}\right)^2 a_n\).

\begin{enumerate}
\item 证明数列 \(\left\{\frac{a_n}{n^2}\right\}\) 是等比数列，并求 \(\{a_n\}\) 的通项公式；

\item 令 \(b_n=a_{n+1}-\frac{1}{2}a_n\)，求数列 \(\{b_n\}\) 的前 \(n\) 项和 \(S_n\)；

\item 求数列 \(\{a_n\}\) 的前 \(n\) 项和 \(T_n\).
\end{enumerate}
\end{problem}

\begin{solution}
\begin{enumerate}
\item 由递推式
\[
2a_{n+1}=\left(\frac{n+1}{n}\right)^2a_n
\]
得
\[
\frac{a_{n+1}}{(n+1)^2}=\frac12\frac{a_n}{n^2}
\]
因此数列 \(\left\{\dfrac{a_n}{n^2}\right\}\) 的相邻两项之比恒为 \(\dfrac12\). 又 \(a_1=1\)，所以其首项为 \(\dfrac{a_1}{1^2}=1\)，从而 \(\frac{a_n}{n^2}=\frac{1}{2^{n-1}}\)，\(a_n=\frac{n^2}{2^{n-1}}\).

\item 由（1）
\[
b_n=a_{n+1}-\frac12a_n=\frac{(n+1)^2}{2^n}-\frac{n^2}{2^n}=\frac{2n+1}{2^n}
\]
所以
\[
S_n=\frac32+\frac5{2^2}+\frac7{2^3}+\cdots+\frac{2n+1}{2^n}
\]
两边乘以 \(\dfrac12\)，得
\[
\frac12S_n=\frac3{2^2}+\frac5{2^3}+\frac7{2^4}+\cdots+\frac{2n-1}{2^n}+\frac{2n+1}{2^{n+1}}
\]
两式相减，得到
\[
\frac12S_n=\frac32+2\left(\frac1{2^2}+\frac1{2^3}+\cdots+\frac1{2^n}\right)-\frac{2n+1}{2^{n+1}}
\]
又 \(\displaystyle\frac1{2^2}+\frac1{2^3}+\cdots+\frac1{2^n}=\frac12-\frac1{2^n}\)，
于是 \(\frac12S_n=\frac52-\frac{2}{2^n}-\frac{2n+1}{2^{n+1}}\)，从而 \(S_n=5-\frac{2n+5}{2^n}\).

\item 由定义
\[
S_n=\sum_{k=1}^n\left(a_{k+1}-\frac12a_k\right)=\left(a_2+a_3+\cdots+a_{n+1}\right)-\frac12\left(a_1+a_2+\cdots+a_n\right)
\]
令 \(T_n=a_1+a_2+\cdots+a_n\)，则
\[
S_n=(T_n-a_1+a_{n+1})-\frac12T_n=\frac12T_n-a_1+a_{n+1}
\]
因此
\[
T_n=2S_n+2a_1-2a_{n+1}=2S_n+2-2\frac{(n+1)^2}{2^n}=12-\frac{n^2+4n+6}{2^{n-1}}
\]
\end{enumerate}
\end{solution}

\begin{problem}
已知椭圆 \(C_1\) 的中心和抛物线 \(C_2\) 的顶点都在坐标原点 \(O\)，\(C_1\) 和 \(C_2\) 有公共焦点 \(F\)，点 \(F\) 在 \(x\) 轴上，且在 \(C_1\) 的长轴上，且 \(C_1\) 的长轴长，短轴长及点 \(F\) 到 \(C_1\) 右准线的距离成等比数列.

\begin{enumerate}
\item 当 \(C_2\) 的准线与 \(C_1\) 右准线间的距离为 \(15\) 时，求 \(C_1\) 及 \(C_2\) 的方程；

\item 设过点 \(F\) 且斜率为 \(1\) 的直线 \(l\) 交 \(C_1\) 于 \(P\)，\(Q\) 两点，交 \(C_2\) 于 \(M\)，\(N\) 两点. 当 \(|MN|=8\) 时，求 \(|PQ|\) 的值.
\end{enumerate}
\end{problem}

\begin{solution}
\begin{enumerate}
\item 设 \(C_1:\frac{x^2}{a^2}+\frac{y^2}{b^2}=1\)（\(a>b>0\)），其半焦距为 \(c\)（\(c>0\)）. 则 \(C_2:y^2=4cx\). 由条件知 \((2b)^2=2a\left(\frac{a^2}{c}-c\right)\)，得 \(a=2c\). \(C_1\) 的右准线方程为 \(x=\frac{a^2}{c}\)，即 \(x=4c\). \(C_2\) 的准线方程为 \(x=-c\). 由条件知 \(5c=15\)，所以 \(c=3\)，故 \(a=6\)，\(b=3\sqrt{3}\). 从而 \(C_1:\frac{x^2}{36}+\frac{y^2}{27}=1\)，\(C_2:y^2=12x\).

\item 由初始条件，设椭圆半焦距为 \(c\). 由第（1）问中的关系仍有 \(a=2c\)、\(b=\sqrt3c\)，故 \(C_1:3x^2+4y^2=12c^2\)，\(C_2:y^2=4cx\).
焦点为 \(F(c,0)\)，过 \(F\) 且斜率为 \(1\) 的直线为 \(l:y=x-c\).

先求椭圆弦 \(PQ\). 将 \(y=x-c\) 代入椭圆方程，得
\[
3x^2+4(x-c)^2=12c^2
\]
即
\[
7x^2-8cx-8c^2=0
\]
设其两个根为 \(x_3\)，\(x_4\)，由求根公式，或由根的判别式，
\[
|x_3-x_4|=\frac{\sqrt{(-8c)^2-4\cdot7\cdot(-8c^2)}}{7}
=\frac{12\sqrt2}{7}c
\]
直线 \(l\) 的斜率为 \(1\)，所以两交点的纵坐标差等于横坐标差，从而
\[
|PQ|=\sqrt{(x_3-x_4)^2+(y_3-y_4)^2}
=\sqrt2|x_3-x_4|=\frac{24}{7}c
\]

再求抛物线弦 \(MN\). 将 \(y=x-c\) 代入 \(C_2\)，得
\[
(x-c)^2=4cx
\]
即
\[
x^2-6cx+c^2=0
\]
设其两个根为 \(x_1\)，\(x_2\). 则
\[
|x_1-x_2|=\sqrt{(-6c)^2-4c^2}=4\sqrt2c
\]
同样因为直线斜率为 \(1\)，
\[
|MN|=\sqrt2|x_1-x_2|=8c
\]
题设 \(|MN|=8\)，故 \(c=1\). 因此
\[
|PQ|=\frac{24}{7}c=\boxed{\frac{24}{7}}
\]
\end{enumerate}
\end{solution}

\begin{problem}
设函数 \(f(x)=x^3-x^2-x+2\).

\begin{enumerate}
\item 求 \(f(x)\) 的单调区间和极值；

\item 若当 \(x\in[-1,2]\) 时，\(-3\le af(x)+b\le 3\)，求 \(a-b\) 的最大值.
\end{enumerate}
\end{problem}

\begin{solution}
\begin{enumerate}
\item \(f'(x)=3x^2-2x-1=(3x+1)(x-1)\). 于是当 \(x\in\left(-\frac{1}{3},1\right)\) 时，\(f'(x)<0\)；当 \(x\in\left(-\infty,-\frac{1}{3}\right)\cup(1,+\infty)\) 时，\(f'(x)>0\). 故 \(f(x)\) 在 \(\left(-\frac{1}{3},1\right)\) 单调减少，在 \(\left(-\infty,-\frac{1}{3}\right)\cup(1,+\infty)\) 单调增加. 当 \(x=-\frac{1}{3}\) 时，\(f(x)\) 取得极大值 \(f\left(-\frac{1}{3}\right)=\frac{59}{27}\). 当 \(x=1\) 时，\(f(x)\) 取得极小值 \(f(1)=1\).

\item 根据（1）及 \(f(1)=1\)，\(f(2)=4\)，可知 \(f(x)\) 在 \([-1,2]\) 的最大值为 \(4\)，最小值为 \(1\). 因此，当 \(x\in[-1,2]\) 时，\(-3\le af(x)+b\le 3\) 的充要条件是 \(-3\le a+b\le 3\) 且 \(-3\le 4a+b\le 3\)，即 \(a\)，\(b\) 满足约束条件
\[
\begin{cases}
a+b\ge -3,\\
a+b\le 3,\\
4a+b\ge -3,\\
4a+b\le 3.
\end{cases}
\]
令 \(u=a+b\)，\(v=4a+b\).
约束条件等价于 \(-3\le u\le3\)、\(-3\le v\le3\). 由 \(v-u=3a\)，\(4u-v=3b\)，
可得
\[
a-b=\frac{v-u}{3}-\frac{4u-v}{3}=\frac{2v-5u}{3}
\]
因为 \(v\le3\) 且 \(u\ge-3\)，所以
\[
a-b=\frac{2v-5u}{3}\le\frac{2\cdot3-5\cdot(-3)}{3}=7
\]
当 \(u=-3\)，\(v=3\) 时，\(a=\dfrac{v-u}{3}=2\)，\(b=\dfrac{4u-v}{3}=-5\)，满足全部约束，且 \(a-b=7\). 因此 \(a-b\) 的最大值为 \(7\).
\end{enumerate}
\end{solution}
